\begin{table}[htbp]
  \centering
  \caption{Cost against element count on a repeatedly refined surface, each level quadrupling the element count. Every time is one collision step at that level, in milliseconds. \emph{BP full} is the whole broad phase, the acceleration structure plus the queries over it; \emph{BP prep} is the structure alone and \emph{BP queries} the queries alone. The narrow phase is named by the answer it is asked for: \emph{NP EToI} returns one earliest time of impact for the step, so every query prunes against the running minimum, and \emph{NP per-pair} returns one per candidate with no shared bound. BP prep here is the acceleration structure and BP queries the two overlap queries over it. $p$ is the least-squares exponent in $t \sim n^{p}$ fitted over all levels of that series.}
  \label{tab:scaling}
  \fittable{%
\begin{tabular}{lrrrrrrrr}
    \toprule
    series & level & elements & candidate pairs & BP prep (ms) & BP queries (ms) & BP full (ms) & NP EToI (ms) & p \\
    \midrule
    CPU / cell2d & 0 & 92,230 & 17,982 & 0.9 & 3.1 & 4.0 & 0.4 & 1.00 \\
     & 1 & 368,920 & 87,771 & 4.7 & 14.2 & 18.9 & 0.6 &  \\
     & 2 & 1,475,680 & 379,818 & 14.1 & 53.8 & 67.9 & 2.1 &  \\
     & 3 & 5,902,720 & 1,573,741 & 52.4 & 204.7 & 257.0 & 7.7 &  \\
     & 4 & 23,610,880 & 6,402,081 & 233.3 & 765.3 & 998.7 & 30.3 &  \\
     & 5 & 94,443,520 & 25,824,702 & 1195.1 & 3596.4 & 4791.5 & 118.2 &  \\
    CPU / sweep & 0 & 92,230 & 17,982 & 2.1 & 2.6 & 4.7 & 0.2 & 1.43 \\
     & 1 & 368,920 & 87,771 & 8.1 & 16.5 & 24.6 & 0.7 &  \\
     & 2 & 1,475,680 & 379,818 & 26.9 & 146.6 & 173.5 & 2.3 &  \\
     & 3 & 5,902,720 & 1,573,741 & 105.2 & 1230.4 & 1335.6 & 9.4 &  \\
     & 4 & 23,610,880 & 6,402,081 & 418.8 & 11052.0 & 11470.7 & 37.2 &  \\
     & 5 & 94,443,520 & 25,824,702 & 2172.9 & 89521.1 & 91694.1 & 146.2 &  \\
    GPU / cell2d & 0 & 92,230 & 17,982 & 0.7 & 1.8 & 2.5 & 1.6 & 1.18 \\
     & 1 & 368,920 & 87,771 & 16.1 & 4.5 & 20.6 & 2.0 &  \\
     & 2 & 1,475,680 & 379,818 & 61.3 & 11.0 & 72.2 & 1.9 &  \\
     & 3 & 5,902,720 & 1,573,741 & 244.0 & 40.2 & 284.2 & 2.6 &  \\
     & 4 & 23,610,880 & 6,402,081 & 1482.8 & 342.3 & 1825.1 & 5.6 &  \\
     & 5 & 94,443,520 & 25,824,702 & 11724.2 & 9162.4 & 20886.6 & 7.6 &  \\
    GPU / sweep & 0 & 92,230 & 17,982 & 4.4 & 3.9 & 8.3 & 1.9 & 1.19 \\
     & 1 & 368,920 & 87,771 & 29.9 & 11.6 & 41.4 & 1.8 &  \\
     & 2 & 1,475,680 & 379,818 & 111.3 & 50.6 & 161.9 & 2.4 &  \\
     & 3 & 5,902,720 & 1,573,741 & 452.5 & 354.9 & 807.5 & 5.4 &  \\
     & 4 & 23,610,880 & 6,402,081 & 1962.6 & 3392.7 & 5355.3 & 6.6 &  \\
     & 5 & 94,443,520 & 25,824,702 & 15523.3 & 26705.0 & 42228.3 & 38.9 &  \\
    \bottomrule
  \end{tabular}}
  \par\smallskip\footnotesize The two frames used here do not come into contact, so the narrow phase has almost no work to do and its column is dominated by noise rather than by element count; what this measures is the broad phase and the preparation that feeds it. Narrow-phase cost against problem size is in the per-case figure, over cases that do collide. Where the exponent is below 1 it is because the fixed cost visible at the smallest size is amortised as the mesh grows. \texttt{prep} builds the acceleration structure -- the cell list's grid or the sweep's sorted intervals -- and \texttt{step} is the traversal that reports pairs; the two strategies divide the work between those columns quite differently.
  \par\smallskip\footnotesize Source: \texttt{benchmark/results/scaling/}
\end{table}
